\documentclass[12pt, oneside]{article}

\usepackage[letterpaper, scale=0.89, centering]{geometry}
\usepackage{fancyhdr}
\setlength{\parindent}{0em}
\setlength{\parskip}{1em}

\pagestyle{fancy}
\fancyhf{}
\renewcommand{\headrulewidth}{0pt}
\rfoot{\href{https://creativecommons.org/licenses/by-nc-sa/2.0/}{CC BY-NC-SA 2.0} Version \today~(\thepage)}

\usepackage{amssymb,amsmath,pifont,amsfonts,comment,enumerate,enumitem}
\usepackage{currfile,xstring,hyperref,tabularx,graphicx,wasysym}
\usepackage[labelformat=empty]{caption}
\usepackage[dvipsnames,table]{xcolor}
\usepackage{multicol,multirow,array,listings,tabularx,lastpage,textcomp,booktabs}

\lstnewenvironment{algorithm}[1][] {   
    \lstset{ mathescape=true,
        frame=tB,
        numbers=left, 
        numberstyle=\tiny,
        basicstyle=\rmfamily\scriptsize, 
        keywordstyle=\color{black}\bfseries,
        keywords={,procedure, div, for, to, input, output, return, datatype, function, in, if, else, foreach, while, begin, end, }
        numbers=left,
        xleftmargin=.04\textwidth,
        #1
    }
}
{}
\lstnewenvironment{java}[1][]
{   
    \lstset{
        language=java,
        mathescape=true,
        frame=tB,
        numbers=left, 
        numberstyle=\tiny,
        basicstyle=\ttfamily\scriptsize, 
        keywordstyle=\color{black}\bfseries,
        keywords={, int, double, for, return, if, else, while, }
        numbers=left,
        xleftmargin=.04\textwidth,
        #1
    }
}
{}

\newcommand\abs[1]{\lvert~#1~\rvert}
\newcommand{\st}{\mid}

\newcommand{\A}[0]{\texttt{A}}
\newcommand{\C}[0]{\texttt{C}}
\newcommand{\G}[0]{\texttt{G}}
\newcommand{\U}[0]{\texttt{U}}

\newcommand{\cmark}{\ding{51}}
\newcommand{\xmark}{\ding{55}}

 
\begin{document}
\begin{flushright}
    \StrBefore{\currfilename}{.}
\end{flushright} 
\section*{Monday November 15}



We have the following subset relationships between sets of numbers:

\[
    \mathbb{Z}^{+} \subsetneq \mathbb{N} \subsetneq \mathbb{Z} \subsetneq \mathbb{Q} \subsetneq \mathbb{R}
\]


Which of the proper subset inclusions above can you prove?

\vspace{50pt} 

{\bf Definition}: A {\bf finite} set is one whose distinct elements can be counted by a natural number.
 

{\bf Motivating question}: when can we say one set is {\it bigger than} another?

Which is bigger? 
\begin{itemize}
    \item The set $\{1,2,3\}$ or the set $\{0,1,2,3\}$?
    \item The set $\{0, \pi, \sqrt{2} \}$ or the set $\{\mathbb{N}, \mathbb{R}, \emptyset\}$?
    \item The set $\mathbb{N}$ or the set $\mathbb{R}^+$?
\end{itemize}

{\it Which of the sets above are finite? infinite?} 

{\bf Key idea for cardinality}: Counting 
distinct elements is a way of labelling elements
with natural numbers. This is a function!
In general, functions let us 
associate elements of one set with those
of another. If the association is ``{\it good}", 
we get a correspondence between the elements of the subsets
which can relate the sizes of the sets. 

{\it Analogy}: Musical chairs

\begin{multicols}{2}
\includegraphics[width=1.8in]{../../resources/images/musicalchairs.png}
\columnbreak

People try to sit down when the music stops

Person\sun~ sits in Chair 1,
Person\smiley~ sits in Chair 2,

Person\frownie~  is left standing!
\end{multicols}
What does this say about the number of chairs and the number of people?

\vspace{100pt} \newpage


Recall that a function is defined by its (1) domain, (2) codomain, and (3) rule assigning each 
element in the domain exactly one element in the codomain. 
The domain and codomain are nonempty sets.
The rule can be depicted as a table, formula, English description, etc.

A function can {\it fail to be well-defined} if there is some 
domain element where the function rule doesn't give a
unique codomain element. For example, the function rule might lead to 
more than one potential image, or to an image outside of the codomain.


{\it Example}: $f_A: \mathbb{R}^+ \to \mathbb{Q}$ with $f_A(x) = x$ is {\bf not} a well-defined function because

\vspace{100pt}


{\it Example}: $f_B: \mathbb{Q} \to \mathbb{Z}$ with $f_B\left(\frac{p}{q}\right) = p+q$ is {\bf not} a well-defined function because

\vspace{100pt}


{\it Example}: $f_C: \mathbb{Z} \to \mathbb{R}$ with $f_C(x) = \frac{x}{|x|}$ is {\bf not} a well-defined function because

\vspace{100pt}
 

{\bf Definition} : A function $f: D  \to C$ is {\bf one-to-one} (or  injective) 
means for every $a,b$ in the domain $D$, 
if $f(a) = f(b)$ then  $a=b$.

Formally, $f: D  \to  C$ is  one-to-one  means $\underline{\phantom{\forall a \in D \forall b \in D ~(f(a) = f(b) \to a = b)}}$.
 

Informally, a function being one-to-one means ``no duplicate images''.

\phantom{Draw finite domain, finite codomain picture with duplicate image.}
\vspace{50pt} 

{\bf Definition}:  For nonempty sets $A, B$, we say that {\bf the  cardinality of $A$ is  no  bigger than the cardinality of $B$}, 
and write $|A| \leq |B|$, to mean there is a  one-to-one function  with domain $A$  and codomain $B$.
Also, we define $|\emptyset| \leq |B|$ for all sets $B$. 

{\it In the analogy}: The function $sitter: \{ Chair1, Chair2\} \to \{ Person\sun, Person\smiley, Person\frownie \}$ given
by $sitter(Chair1) = Person\sun$,  $sitter(Chair2) = Person\smiley$, is one-to-one and witnesses that 
\[
| \{ Chair1, Chair2\} | \leq |\{ Person\sun, Person\smiley, Person\frownie \}|
\] 

Let $S_2$ be the set of RNA strands of length 2, formally $S_2 = \{ s \in S \mid rnalen(s) = 2\}$.

{\bf True} or {\bf False}: $| \{\A,\U,\G,\C\} | \leq |S_2 |$ 

{\it Why?}
\vspace{80pt}

{\bf True} or {\bf False}: $| \{\A,\U,\G,\C\} \times \{\A, \U, \G,\C\} | \leq |S_2 |$

{\it Why?}
\vspace{80pt}
 

{\bf Definition}: A function $f: D  \to C$ is {\bf onto} (or  surjective) means for every $b$ in the codomain, 
there  is an element $a$ in the domain with  $f(a) = b$.

Formally, $f: D  \to  C$ is  onto  means $\underline{\phantom{\forall b \in C  \exists a \in D ( f(a) = b)}}$.
 

Informally, a function being onto means ``every potential image is an actual image''.

\phantom{Draw finite domain, finite codomain picture with duplicate image.}
\vspace{50pt} 

{\bf Definition}:  For nonempty sets $A, B$, we say that {\bf the  cardinality of $A$ is  no  smaller than 
the cardinality of  $B$}, and 
write $|A| \geq |B|$, to mean there is an onto function  with domain $A$  and codomain $B$.
Also, we define $|A| \geq |\emptyset|$ for all sets $A$.
 

{\it In the analogy}: The function $triedToSit: \{ Person\sun, Person\smiley, Person\frownie \} \to  \{ Chair1, Chair2\} $ given
by $triedToSit(Person\sun) = Chair1$,  $triedToSit(Person\smiley) = Chair2$, 
$triedToSit(Person\frownie) = Chair2$, is onto and witnesses that 
\[
 |\{ Person\sun, Person\smiley, Person\frownie \}| \geq | \{ Chair1, Chair2\} |
\] \newpage


Let $S_2$ be the set of RNA strands of length 2.

{\bf True} or {\bf False}: $ |S_2 | \geq | \{\A,\U,\G,\C\} |$

{\it Why?}
\vspace{80pt}

{\bf True} or {\bf False}: $ |S_2 | \geq | \{\A,\U,\G,\C\} \times \{\A, \U, \G,\C\} |$

{\it Why?}
\vspace{80pt} 

{\bf Definition} : A function $f: D  \to C$ is a {\bf bijection} means that it is both 
one-to-one  and onto. The {\bf inverse} of a  bijection $f: D  \to  C$ is 
the function $g: C  \to  D$  such that $g(b) = a$ iff  $f(a) =  b$.
 \subsection*{Cardinality of sets}

For nonempty sets $A, B$ we say
\begin{align*}
|A| \leq |B| &\text{ means there is a one-to-one function with domain $A$, codomain $B$} \\
|A| \geq |B| &\text{ means there is an onto function with domain $A$, codomain $B$} \\
|A| = |B| &\text{ means there is a bijection with domain $A$, codomain $B$}
\end{align*}

For all sets $A$, we say $|A| = |\emptyset|$, $|\emptyset| = |A|$ if and only if $A = \emptyset$. 

{\it Caution}: we use 
familiar symbols to define cardinality, like
$| \phantom{\cdots} | \leq | \phantom{\cdots} |$
and 
$| \phantom{\cdots} | \geq | \phantom{\cdots} |$
and 
$| \phantom{\cdots} | = | \phantom{\cdots} |$, 
but the meaning of these symbols depends on context.
We've seen that vertical lines can mean absolute
value (for real numbers), divisibility (for integers), 
and now sizes (for sets). 

Now we see that $\leq$ and $\geq$ can mean comparing
numbers or comparing sizes of sets. When the sets being 
compared are finite, the definitions 
of $|A| \leq |B|$ agree. 

But, properties of numbers cannot be assumed when comparing
cardinalities of infinite sets.

In a nutshell: cardinality of sets is defined via functions.
This definition agrees with the usual notion of ``size'' for 
finite sets. \newpage
\section*{Review}
\begin{enumerate}
    \item \hspace{1in}\\ 

Select all and only the {\bf finite} sets below.

\begin{enumerate}
\item $X = \{ a,b,c\}$
\item $Y = \{ 1, 2, 3, 4, 5\}$
\item $Z = \{ 10, 20, 30 \}$
\item $\emptyset$
\item $\mathbb{Z}$
\item $\{ \emptyset \}$
\item $\{ \mathbb{Z} \}$
\end{enumerate}     \item \hspace{1in}\\ Consider the following input-output definition tables with $X = \{ a,b,c\}$ and 
$Y = \{ 1, 2, 3, 4, 5\}$ and $Z = \{ 10, 20, 30 \}$

\begin{center}
\begin{multicols}{3}
\begin{tabular}{c}
Table 1\\
 \begin{tabular}{c|c}
Input & Output \\
\hline
$1$ & $10$ \\
$2$ & $20$ \\
$3$ & $30$ \\
\end{tabular}\end{tabular}
\columnbreak
\begin{tabular}{c}
Table 2\\ 
\begin{tabular}{c|c}
Input & Output \\
\hline
$a$ & $1$ \\
$b$ & $4$ \\
$c$ & $5$ \\
\end{tabular}\end{tabular}
\columnbreak
\begin{tabular}{c}
Table 3\\ 
\begin{tabular}{c|c}
Input & Output \\
\hline
$10$ & $a$ \\
$20$ & $b$ \\
$30$ & $a$ \\
\end{tabular}\end{tabular}
\end{multicols}
\end{center}

\begin{enumerate}
\item Select all and only the tables that each define a well-defined function 
whose domain and codomain 
is each $X$, $Y$, or $Z$.
\item Select all and only the tables that each define a well-defined function 
(with domain $X$ or $Y$ or $Z$ and 
with codomain $X$ or $Y$ or $Z$) and that is one-to-one.
\item Select all and only the tables that each define a well-defined function 
(with domain $X$ or $Y$ or $Z$ and 
with codomain $X$ or $Y$ or $Z$) and that is onto.
\end{enumerate}     \newpage
    \item \hspace{1in}\\ 

Consider the following functions:

\[
\begin{array}{l|l}
\begin{array}{ll}
f : \mathbb{Z} &\to \mathbb{N} \\
f(n) & = \begin{cases}
  0 & \textrm{ when } n = 0 \\
  (-2 \cdot n) - 1 & \textrm{ when } n < 0 \\
  2 \cdot n & \textrm{ when } n > 0
\end{cases}
\end{array}
&
\begin{array}{ll}
g : \mathbb{Z} &\to \mathbb{N} \\
g(n) & = \begin{cases}
  -1 \cdot n & \textrm{ when } n < 0 \\
  n & \textrm{ when } n \geq 0
\end{cases}
\end{array} \\
\hline

\begin{array}{ll}
h : \mathbb{N} &\to \mathbb{Z} \\
h(n) & = \begin{cases}
  (-2 \cdot n) + 1 & \textrm{ when } n \textrm{ is even } \\
  2 \cdot n & \textrm{ when } n \textrm{ is odd }
\end{cases}
\end{array}
&
\begin{array}{ll}
q : \mathbb{N} &\to \mathbb{Z} \\
q(n) & = \begin{cases}
  -1 \cdot ((n + 1) \textbf{ div } 2) & \textrm{ when } n \textrm{ is odd} \\
  n \textbf{ div } 2 & \textrm{ when } n \textrm{ is even} \\
\end{cases}
\end{array} \\
\end{array}
\]

\begin{enumerate}
\item What is the result of $f(-3)$?

\item What is the result of $q(f(-4))$?

 {\it Notice we are looking at function composition here: first apply $f$ and then apply 
 $q$ to the result.}

\item What is the result of $f(h(4))$? 

{\it Notice we are looking at function composition here: first apply $h$ and then apply 
$f$ to the result.}

\item What is the result of $g(-4)$?

\item What is the result of $g(4)$?

\item Consider the following statements, and indicate if they are true for each of $f$, $g$, $h$, and $q$.

\begin{enumerate}[label=\roman*.]
    \item This function is one-to-one.
    \item This function is onto.
    \item This function is a bijection.
    \item This function could serve as a witness for $|\mathbb{Z}| \leq |\mathbb{N}|$.
    \item This function could serve as a witness for $|\mathbb{Z}| \geq |\mathbb{N}|$.
    \item This function could serve as a witness for $|\mathbb{N}| \leq |\mathbb{Z}|$.
    \item This function could serve as a witness for $|\mathbb{N}| \geq |\mathbb{Z}|$.
\end{enumerate}

\end{enumerate} \end{enumerate}
\newpage

\section*{Wednesday November 17}


{\bf Properties of cardinality}
\begin{align*}
&\forall A ~ (~  |A| = |A| ~)\\
&\forall A ~ \forall B ~(~ |A| = |B|  ~\to ~ |B| = |A|~)\\
&\forall A ~ \forall B ~ \forall C~ (~ (|A| = |B| ~\wedge~ |B| = |C|) ~\to ~ |A| = |C|~)
\end{align*}

{\it Extra practice with proofs:} Use the definitions of bijections to prove these properties. \vspace{100pt}


{\bf Cantor-Schroder-Bernstein Theorem}: For all nonempty sets,
\[
|A| = |B| \qquad\text{if and only if} \qquad (|A| \leq |B| ~\text{and}~ |B| \leq |A|)
\qquad\text{if and only if} \qquad (|A| \geq |B| ~\text{and}~ |B| \geq |A|)
\]

To prove $|A| = |B|$,  we can do any {\bf one} of the following

\begin{itemize}\setlength{\itemsep}{-5pt}
\item Prove there exists  a bijection $f:  A \to B$;
\item Prove there exists a  bijection  $f: B  \to  A$;
\item Prove there exists two functions $f_1: A \to B$, $f_2: B \to  A$ where each of $f_1, f_2$ is one-to-one.
\item Prove there exists two functions $f_1: A \to B$, $f_2: B \to  A$ where each of $f_1, f_2$ is onto.
\end{itemize} \newpage


{\bf Definition}: A set $A$ is {\bf countably infinite} means it is the 
same size as $\mathbb{N}$.

 

{\bf Natural numbers} $\mathbb{N}$ \hfill {\it List}:  $0~~1~~2~~3~~4~~5~~6~~7~~8~~9~~10 \ldots$

$identity: \mathbb{N} \to \mathbb{N}$ with $identity(n) = n$

{\it Claim}: $identity$ is a bijection. {\it Proof}: Ex. \hfill {\bf Corollary}: $ | \mathbb{N} | = |\mathbb{N}|~$

{\bf Positive integers} $\mathbb{Z}^+$ \hfill {\it List}:  $1~~2~~3~~4~~5~~6~~7~~8~~9~~10~~11\ldots$

$positives: \mathbb{N} \to \mathbb{Z}^+$ with $positives(n) = n+1$

{\it Claim}: $positives$ is a bijection.  {\it Proof}: Ex.\hfill {\bf Corollary}: $ | \mathbb{N} | = |\mathbb{Z}^+|$

{\bf Negative integers $\mathbb{Z}^-$}\hfill  {\it List}: $-1$~$-2$~$-3$~$-4$~$-5$~$-6$~$-7$~$-8$~$-9$~$-10$~$-11$\ldots

$negatives: \mathbb{N} \to \mathbb{Z}^-$ with $negatives(n) = -n-1$

{\it Claim}: $negatives$ is a bijection. \hfill {\bf Corollary}: $ | \mathbb{N} | = |\mathbb{Z}^-|$

{\it Proof}: We need to show it is a well-defined function that is one-to-one and onto.

\begin{itemize}
\item Well-defined? 

Consider an arbitrary element of the domain, $n \in \mathbb{N}$. We need to show it maps to exactly one element of $\mathbb{Z}^-$.

\vfill

\item One-to-one?


Consider arbitrary elements of the domain $a, b \in \mathbb{N}$. We need to show that 
$$\left(~negatives(a) = negatives(b) ~\right) \to (a=b)$$

\vfill

\item Onto?

Consider arbitrary element of the codomain $b \in \mathbb{Z}^-$. We need witness in $\mathbb{N}$ that maps to $b$.

\vfill
\end{itemize}

{\bf Integers} $\mathbb{Z}$ \hfill {\it List}:  $0~-1~~1~-2~~2~-3~~3~-4~~4~-5~~5$\ldots

$f: \mathbb{Z} \to \mathbb{N}$ with $f(x) = \begin{cases}2x &\text{if $x \geq 0$} \\-2x-1 &\text{if $x < 0$} \end{cases}$

{\it Claim}: $f$ is a bijection.  {\it Proof}: Ex.\hfill {\bf Corollary}: $ | \mathbb{Z} | = |\mathbb{N}|$
\newpage


 

{\bf More examples of countably infinite sets}

{\bf Claim}: $S$ is countably infinite

{\it Similarly: The set of all strings over a specific alphabet is countably infinite.}
\begin{center}
Bijection using alphabetical-ish ordering
(first order by length, then alphabetically among strings of same length) 
of strands
\end{center}

\vspace{50pt}


{\bf Claim}: $L$ is countably infinite

\begin{multicols}{2}
\begin{align*}
    &list: \mathbb{N} \to L \\
    &list(n) = (n, []) \\
    &
\end{align*}

\begin{align*}
    &toNum: L \to \mathbb{N} \\
    &toNum([]) = 0 \\
    &toNum( ~(n,l)~) = 2^n 3^{toNum}(l) \qquad \text{for $n \in \mathbb{N}$, $l \in L$}
\end{align*}
\end{multicols}
\vspace{30pt}

{\bf Claim}: $|\mathbb{Z}^+| = |\mathbb{Q}|$ 

One-to-one function from $\mathbb{Z}^+$ to $\mathbb{Q}$
is $f_1: \mathbb{Z} \to \mathbb{Q}$ with $f_1(n) = n$ for all 
$n \in \mathbb{N}$.

\vspace{30pt}


\begin{align*}
    &f_2: \mathbb{Q} \to \mathbb{Z} \times \mathbb{Z} \\
    &f_2(x) = \begin{cases}
        (0,1) & \text{if $x = 0$} \\
        (p,q) & \text{if $x = \frac{p}{q}$,}\\
                & \text{$gcd(p,q) = 1$, $q > 0$}
    \end{cases}
\end{align*}
\begin{multicols}{2}
\begin{align*}
    &f_3: \mathbb{Z} \times \mathbb{Z} \to \mathbb{Z}^+ \times \mathbb{Z}^+ \\
    &f_3(~(x,y)~) = \begin{cases}
        (2x+2, 2y+2) & \text{if $x \geq 0, y \geq 0$} \\
        (-2x-1, 2y+2) & \text{if $x < 0, y \geq 0$} \\
        (2x+2, -2y+1) & \text{if $x \geq 0, y < 0$} \\
        (-2x-1, -2y-1) & \text{if $x < 0, y < 0$} \\
    \end{cases}
\end{align*}

\begin{align*}
    &f_4: \mathbb{Z}^+ \times \mathbb{Z}^+ \to \mathbb{Z}^+ \\
    &f_4(~(x,y)~) = 2^x 3^y \qquad \text{for $x,y \in \mathbb{Z}^+$}
\end{align*}
\end{multicols} \newpage
\subsection*{Review}
\begin{enumerate}
    \item 

Consider the function $f: \mathbb{N} \to  \mathbb{Z}$ given by 
$f(n)  =  \begin{cases}
    n~\text{\bf div}~4  \qquad &\text{if $n$ is even} \\
    -( (n+1)~\text{\bf div}~4) \qquad &\text{if $n$ is odd}
\end{cases}$
    
Select all and only the true statements below.
    \begin{enumerate}
        \item $f$ is one-to-one
        \item $f$ is onto
        \item $f$ is a bijection
        \item $f$ witnesses that $|\mathbb{N}| \leq |\mathbb{Z}|$
        \item $f$ witnesses that $|\mathbb{N}| \geq |\mathbb{Z}|$
        \item $f$ witnesses that $|\mathbb{N}| = |\mathbb{Z}|$
        \item There is a one-to-one function
        with domain $\mathbb{N}$ and codomain
        $\mathbb{Z}$
        \item There is an  onto function
        with domain $\mathbb{N}$ and codomain
        $\mathbb{Z}$
        \item There is a bijection
        with domain $\mathbb{N}$ and codomain
        $\mathbb{Z}$
    
        \item $|\mathbb{N}| \leq |\mathbb{Z}|$
        \item $|\mathbb{N}| \geq |\mathbb{Z}|$
        \item $|\mathbb{N}| = |\mathbb{Z}|$
    \end{enumerate}     \item \hspace{1in}\\ 

{\it  Goals for this question: Reason through
multiple nested quantifiers. Fluently use the definition and properties of the set of rationals. 
}

Recall the definition of the set of rational numbers, $\mathbb{Q} = \left\{ \frac{p}{q} \mid p \in \mathbb{Z}  \text{ and  } q  \in \mathbb{Z} \text{ and } q \neq  0 \right\}$.
We define the set of {\bf irrational} numbers $\overline{\mathbb{Q}} = \mathbb{R} - \mathbb{Q}
= \{ x \in \mathbb{R} \mid x \notin \mathbb{Q} \}$.
\begin{multicols}{2}
\begin{enumerate}[label=(\roman*)]
\item $\forall x \in \mathbb{Q} ~\forall y \in \mathbb{Q}~ \exists z \in \mathbb{Q} ~( x + y = z)$
\item $\forall x \in \mathbb{Q} ~\forall y \in \mathbb{Q}~ \exists z \in \mathbb{Q} ~( x + z = y)$
\item $\forall x \in \mathbb{Q} ~\forall y \in \mathbb{Q}~ \exists z \in \mathbb{Q} ~( x \cdot y = z)$
\item $\forall x \in \mathbb{Q} ~\forall y \in \mathbb{Q} ~\exists z \in \mathbb{Q} ~( x \cdot z = y)$
\item $\forall x \in \overline{\mathbb{Q}}~ \forall y \in \overline{\mathbb{Q}}~ \exists z \in \overline{\mathbb{Q}} ~( x + y = z)$
\item $\forall x \in \overline{\mathbb{Q}}~ \forall y \in \overline{\mathbb{Q}}~ \exists z \in \overline{\mathbb{Q}}~( x + z = y)$
\item $\forall x \in \overline{\mathbb{Q}} ~\forall y \in \overline{\mathbb{Q}}~ \exists z \in \overline{\mathbb{Q}} ~( x \cdot y = z)$
\item $\forall x \in \overline{\mathbb{Q}} ~\forall y \in \overline{\mathbb{Q}}~ \exists z \in \overline{\mathbb{Q}}~( x \cdot z = y)$
\end{enumerate}
\end{multicols}

\begin{enumerate}
\item Which of the statements above (if any) could be {\bf disproved} using the counterexample 
$x = \frac{1}{2}$, $y= \frac{3}{4}$?
\item Which of the statements above (if any) could be {\bf disproved} using the counterexample 
$x = \sqrt{4}$, $y= \sqrt{3}$?
\item Which of the statements above (if any) could be {\bf disproved} using the counterexample 
$x = 0$, $y= 3$?
\item Which of the statements above (if any) could be {\bf disproved} using the counterexample 
$x = \sqrt{2}$, $y= 0$?
\item Which of the statements above (if any) could be {\bf disproved} using the counterexample 
$x = \sqrt{2}$, $y= -  \sqrt{2}$?
\end{enumerate}

{\it Hint: we proved in class that $\sqrt{2} \notin \mathbb{Q}$. You may also use the facts
that $\sqrt{3} \notin \mathbb{Q}$ and $-\sqrt{2} \notin \mathbb{Q}$.

Bonus - not to hand in: prove these facts; that is, prove that $\sqrt{3} \notin \mathbb{Q}$ and $-\sqrt{2} \notin \mathbb{Q}$. }
 \end{enumerate}

\newpage
\section*{Friday November 19}
\subsection*{Cardinality categories}


A set $A$ is {\bf finite} means it is empty or it is the same size as $\{ 1, \ldots, n \}$ for some $n \in \mathbb{N}$.

A set $A$ is {\bf countably infinite} means it is the same size as $\mathbb{N}$. {\it Notice: 
all countably infinite sets are the same size as each other.}

A set $A$ is {\bf countable} means it is either finite or countably infinite.

A set $A$ is {\bf uncountable} means it is not countable. 

{\bf Lemmas about countable and uncountable sets}

{\bf Lemma}: If $A$ is a subset of a countable set, then it's countable.

\vspace{80pt}

{\bf Lemma}: If $A$ is a superset of an uncountable set, then it's uncountable.

\vspace{80pt}

{\bf Lemma}: If $A$ and $B$ are countable sets, then $A \cup B$ is countable
and $A \cap B$ is countable.

\vspace{80pt}

{\bf Lemma}: If $A$ and $B$ are countable sets, then $A \times B$ is countable.

{\it Generalize pairing ideas from $\mathbb{Z}^+ \times \mathbb{Z}^+$ to $\mathbb{Z}^+$}

\vspace{50pt}

{\bf Lemma}: If $A$ is a subset of $B$ , to show that $|A| = |B|$, 
it's enough to give one-to-one function from $B$ to $A$ or an onto function 
from $A$ to $B$.

\vspace{80pt}
 \newpage
\subsection*{Are there always *bigger* sets?}


{\it Recall}: When $U$ is a set, $\mathcal{P}(U) = \{ X \mid X \subseteq U\}$

{\it Key idea}: For finite sets, the power set of a set has strictly greater size than the set itself.
Does this extend to infinite sets?

{\bf Definition}: For two sets $A, B$, we use the notation $|A| < |B|$ to denote
$(~|A| \leq |B| ~) \land \lnot (~|A| = |B|)$.

\begin{alignat*}{4}
    &\emptyset = \{ \} \qquad &&\mathcal{P}(\emptyset) = \{ \emptyset \} \qquad &&|\emptyset| < |\mathcal{P}(\emptyset)| \\
    &\{1 \} \qquad &&\mathcal{P}(\{1\}) = \{ \emptyset, \{1\} \} \qquad &&|\{1\}| < |\mathcal{P}(\{1\})| \\
    &\{1,2 \} \qquad &&\mathcal{P}(\{1,2\}) = \{ \emptyset, \{1\}, \{2\}, \{1,2\} \} \qquad &&|\{1,2\}| < |\mathcal{P}(\{1,2\})| \\
\end{alignat*}

{\bf $\mathbb{N}$ and its power set}

Example elements of $\mathbb{N}$ 

\vspace{20pt}

Example elements of $\mathcal{P}(\mathbb{N})$

\vspace{20pt}

{\bf Claim}: $| \mathbb{N} | \leq |\mathcal{P} ( \mathbb{N} ) |$

\vspace{100pt}
\newpage
{\bf Claim}: There is an uncountable set.  Example: $\underline{\phantom{~~~\mathcal{P}(\mathbb{N})~~~}}$

{\bf Proof}:  By definition of countable, since $\underline{\phantom{~~~\mathcal{P}(\mathbb{N})~~~}}$
is not finite, {\bf to show} is $|\mathbb{N}| \neq  |\mathcal{P}(\mathbb{N})|$ .

Rewriting using  the definition of  cardinality, {\bf to show} is

\phantom{$\neg \exists f : \mathbb{N} \to \mathcal{P}(\mathbb{N})  ~~(f \text{ is a bijection})~~$}

\phantom{or equivalently $\forall f : \mathbb{N} \to \mathcal{P}(\mathbb{N})  ~~(f \text{ is not a bijection})~~$}


Towards a proof by  universal generalization,  consider  an arbitrary function $f:  \mathbb{N} \to\mathcal{P}(\mathbb{N})$.

{\bf To show}: $f$ is not a bijection.  It's enough to show that $f$ is not onto.

Rewriting using the definition of  onto, {\bf to show}:
\[
\neg  \forall  B \in  \mathcal{P}(\mathbb{N}) ~\exists a \in \mathbb{N}  ~(~f(a) =  B~)
\]
. By logical  equivalence, we can write this as an existential statement:
\[
\underline{\phantom{\qquad\qquad\exists B \in  \mathcal{P}(\mathbb{N}) ~\forall a \in \mathbb{N}  ~(~f(a) \neq  B~)\qquad\qquad}}
\]
In search of a witness, define the following  collection of nonnegative integers:
\[
D_f = \{ n \in \mathbb{N}  ~\mid~  n \notin f(n)  \}
\]
. By  definition  of power  set, since  all elements  of  $D_f$ are  in  $\mathbb{N}$,   $D_f \in \mathcal{P}(\mathbb{N})$.  It's enough to prove the following Lemma: 

{\bf Lemma}: $\forall a \in \mathbb{N}  ~(~f(a) \neq  D_f~)$.


{\bf Proof  of lemma}: \phantom{Towards universal  generalization, consider an arbitrary  $a \in \mathbb{N}$.
By definition  of set equality, {\bf to show} is  $\exists  x ( \neg  (x \in f(a)~  \leftrightarrow  ~x \in D_f))$.
For a witness, consider $x = a$.  There are two cases:  $a \in  f(a)~\vee~a \notin f(a)$. By definition 
of $D_f$, each guarantees that $f(a) \neq  D_f$.}\\

\vspace{50pt}

By  the Lemma, we  have proved that $f$ is not onto, and since $f$ was arbitrary, there are no onto
functions from $\mathbb{N}$ to $\mathcal{P}(\mathbb{N})$. QED


{\bf Where does $D_f$ come from?} The idea is to build a set that would ``disagree" with 
each of the images of $f$ about some element. 

\begin{center}
\begin{tabular}{c|c|ccccccc}
$n \in \mathbb{N}$ & $f(n) = X_n$ &  Is $0   \in X_n$?   & Is $1 \in X_n$?  &  Is $2 \in X_n$?  &  Is $3 \in X_n$?  &
 Is $4 \in X_n$?  &  \ldots & Is $n \in D_f$?\\
\hline
$0$ & $f(0) = X_0$ & {\bf  Y~/~N}  & Y~/~N & Y~/~N & Y~/~N &Y~/~N & \ldots & {\bf  N~/~Y }\\
$1$ & $f(1) = X_1$ & Y~/~N  & {\bf  Y~/~N} & Y~/~N & Y~/~N & Y~/~N & \ldots & {\bf  N~/~Y }\\
$2$ & $f(2) = X_2$ & Y~/~N  & Y~/~N & {\bf  Y~/~N} & Y~/~N &Y~/~N & \ldots & {\bf  N~/~Y }\\
$3$ & $f(3) = X_3$ & Y~/~N  & Y~/~N & Y~/~N & {\bf  Y~/~N} & Y~/~N & \ldots & {\bf  N~/~Y }\\
$4$ & $f(4) = X_4$ & Y~/~N  & Y~/~N & Y~/~N & Y~/~N &{\bf  Y~/~N} & \ldots & {\bf  N~/~Y }\\
\vdots
\end{tabular}
\end{center} \newpage
\subsection*{Countable vs.\ uncountable: sets of numbers}


{\bf Comparing $\mathbb{Q}$ and $\mathbb{R}$} 


Both $\mathbb{Q}$ and $\mathbb{R}$ have no greatest element.

Both $\mathbb{Q}$ and $\mathbb{R}$ have no least element.

The quantified statement 
\[
    \forall x \forall y (x < y \to \exists z ( x < z < y) )
\]
is true about both $\mathbb{Q}$ and $\mathbb{R}$.

Both $\mathbb{Q}$ and $\mathbb{R}$ are infinite. But, $\mathbb{Q}$ is countably infinite
whereas $\mathbb{R}$ is uncountable.\\


{\bf The set of real numbers}

$\mathbb{Z} \subsetneq \mathbb{Q} \subsetneq \mathbb{R}$


{\bf  Order axioms} (Rosen Appendix 1): 

\begin{center}
\begin{tabular}{p{1.2in}p{4in}}
Reflexivity &  $\forall a \in  \mathbb{R} (a \leq a)$\\
Antisymmetry &  $\forall a \in  \mathbb{R}~\forall b \in \mathbb{R}~(~(a \leq b~ \wedge ~b \leq a) \to (a=b)~)$\\
Transitivity &  $\forall a \in  \mathbb{R}~\forall b \in \mathbb{R}~\forall c \in \mathbb{R}~
(~(a \leq b \wedge b \leq c) ~\to  ~(a \leq c)~)$ \\
Trichotomy & 
$\forall a \in \mathbb{R}~\forall b \in \mathbb{R}~ ( ~(a=b ~\vee~ b > a ~\vee~ a  < b)  $
\end{tabular}
\end{center}


{\bf  Completeness axioms} (Rosen Appendix 1): 


\begin{center}
\begin{tabular}{p{1.4in}p{6in}}
Least upper bound &  Every nonempty set of real numbers that 
is bounded  above has  a  least upper bound  
\\
Nested intervals &  For each sequence  of intervals  $[a_n , b_n]$
where, for each $n$, $a_n < a_{n+1} < b_{n+1} < b_n$, there
is at least one  real number $x$ such that, for all $n$, 
$a_n \leq x \leq b_n$.\\
\end{tabular}
\end{center}

Each real  number $r  \in  \mathbb{R}$ is described by a function to give better and better approximations
\[
x_r: \mathbb{Z}^+ \to \{0,1\}  \qquad  \text{where  $x_r(n ) =  n^{th} $ bit in  binary expansion of $r$}
\]
\begin{center}
\begin{tabular}{|c|c|p{3.9in}|}
\hline
$r$ & Binary expansion & $x_r$ \\
\hline
$0.1$ & $0.00011001 \ldots$ &  $x_{0.1}(n) = \begin{cases} 0&\text{if $n > 1$ and $(n~\text{\bf mod}~4) =2$} \\
0&\text{if $n=1$ or if $n > 1$ and $(n~\text{\bf mod}~4) =3$} \\1&\text{if $n > 1$ and $(n~\text{\bf mod}~4) =0$} \\
1&\text{if $n > 1$ and $(n~\text{\bf mod}~4) =1$} \end{cases}$  \\
&&  \\
\hline
$\sqrt{2} - 1 = 0.4142135 \ldots$  &$0.01101010\ldots$& Use linear approximations
(tangent lines from calculus) to get algorithm for bounding error of successive operations. Define 
$x_{\sqrt{2}-1}(n)$ to be  $n^{th}$ bit in approximation  that has error less than  $2^{-(n+1)}$.
\\
&& \\
\hline
\end{tabular}
\end{center}

\newpage 

{\bf Claim}: $\{  r \in \mathbb{R} ~\mid~ 0 \leq r ~\wedge~ r \leq 1 \}$ is uncountable.

{\it Approach 1}: Mimic proof that $\mathcal{P}(\mathbb{Z}^+)$ is uncountable.


{\bf Proof}:  By definition of countable, since $\{  r \in \mathbb{R} ~\mid~ 0 \leq r ~\wedge~ r \leq 1 \}$
is not finite, {\bf to show} is $|\mathbb{N}| \neq  |\{  r \in \mathbb{R} ~\mid~ 0 \leq r ~\wedge~ r \leq 1 \}|$ .


{\bf To show} is
$\forall f : \mathbb{Z}^+ \to \{  r \in \mathbb{R} ~\mid~ 0 \leq r ~\wedge~ r \leq 1 \}  ~~(f \text{ is not a bijection})~~$.
Towards a proof by  universal generalization, consider  an arbitrary function 
$f:  \mathbb{Z}^+ \to \{  r \in \mathbb{R} ~\mid~ 0 \leq r ~\wedge~ r \leq 1 \}$.
{\bf To show}: $f$ is not a bijection.  It's enough to show that $f$ is not onto.
Rewriting using the definition of  onto, {\bf to show}:
\[
\exists x \in \{  r \in \mathbb{R} ~\mid~ 0 \leq r ~\wedge~ r \leq 1 \} ~\forall a \in \mathbb{N}  ~(~f(a) \neq  x~)
\]
In search of a witness, define the following  real number by defining its binary expansion
\[
d_f = 0.b_1 b_2 b_3 \cdots
\]
where $b_i = 1-b_{ii}$ where $b_{jk}$ is the coefficient of $2^{-k}$ in the binary expansion of $f(j)$.
Since\footnote{There's a subtle imprecision in this part of the proof as presented, but it can be fixed.} $d_f \neq f(a)$ for any positive integer $a$, $f$ is not onto.


{\it Approach 2}: Nested closed interval property

{\bf To show} $f: \mathbb{N} \to \{  r \in \mathbb{R} ~\mid~ 0  \leq r ~\wedge~ r \leq 1 \}$ is not onto. 
{\bf  Strategy}: Build
a sequence of nested closed intervals that each avoid some $f(n)$.   Then  the real
number that is in all of the intervals  can't be $f(n)$ for any $n$. Hence,  $f$ is not  onto.

Consider the function $f: \mathbb{N} \to \{  r \in \mathbb{R} ~\mid~ 0 \leq r ~\wedge~ r \leq 1 \}$ with  $f(n) = \frac{1+\sin(n)}{2}$

\begin{center}
\begin{tabular}{c||p{1.65in} || p{3in} }
$n$ &  $f(n)$& Interval that avoids $f(n)$ \\
\hline
$0$ & $0.5$ &  \\
$1$ &$0.920735\ldots$  &  \\
$2$ &$0.954649\ldots$ &  \\
$3$ &$0.570560\ldots$ & \\
$4$ &$0.121599\ldots $&  \\
\vdots &  &\\
\end{tabular}
\end{center}
  \subsection*{Other examples of uncountable sets}


\begin{itemize}
    \item The power set of any countably infinite set is uncountable. For example:
    \[
        \mathcal{P}(\mathbb{N}), \mathcal{P}(\mathbb{Z^+}), \mathcal{P}(\mathbb{Z})
    \]
    are each uncountable.
    \item The closed interval $\{x \in \mathbb{R} ~|~ 0 \leq x \leq 1\}$, any other nonempty closed interval of real numbers whose endpoints are 
    unequal, as well as the related intervals that exclude one or both of the endpoints.
    \item The set of all real numbers $\mathbb{R}$ is uncountable and the set of irrational
    real numbers $\overline{\mathbb{Q}}$ is uncountable.
\end{itemize} \newpage
\subsection*{Review}
\begin{enumerate}
    \item \hspace{1in}\\ 

The diagonalization argument constructs, 
for each function $f:  \mathbb{N} \to  \mathcal{P}(\mathbb{N})$, a set  $D_f$ defined
as
\[
D_f = \{ x \in \mathbb{N} \mid x \notin  f(x) \}
\]
which has the property  that,  for all  $n \in \mathbb{N}$, $f(n) \neq  D_f$.
Consider the following two functions with  domain $\mathbb{N}$ and codomain $\mathcal{P}(\mathbb{N})$
\[
f_1(x) =  \{  y \in  \mathbb{N} \mid y~\text{\bf mod}~3 = x~\text{\bf mod}~3  \}
\]
\[
f_2(x) =  \{  y \in  \mathbb{N} \mid (y > 0) \land
(x ~\text{\bf mod}~y \neq  0)\}
\]

Select all and only the true statements below.
\begin{enumerate}
    \item $0 \in D_{f_1}$
    \item $D_{f_1}$ is infinite
    \item $D_{f_1}$ is uncountable
    \item $1 \in D_{f_2}$
    \item $D_{f_2}$ is empty
    \item $D_{f_2}$ is countably infinite
\end{enumerate}     \item \hspace{1in}\\ 

Recall the definitions from previous assignments and class: 
The bases of RNA are elements of the set 
$B  = \{\A, \C, \G, \U \}$. The set of RNA strands $S$ is defined (recursively) by:

\[
\begin{array}{ll}
\textrm{Basis Step: } & \A \in S, \C \in S, \U \in S, \G \in S \\
\textrm{Recursive Step: } & \textrm{If } s \in S\textrm{ and }b \in B \textrm{, then }sb \in S
\end{array}
\]

For $b$ an integer greater than $1$ and $n$ a positive integer, 
the {\bf base $b$ expansion of $n$}  is
\[
(a_{k-1} \cdots a_1 a_0)_b
\]
where $k$ is a positive integer, $a_0, a_1, \ldots, a_{k-1}$ are nonnegative integers less than $b$, $a_{k-1} \neq  0$, and
\[
n =  a_{k-1} b^{k-1} + \cdots + a_1b + a_0
\]

For $b$ an integer greater than $1$, $w$ a positive integer, and $n$ a nonnegative integer
with $n < b^w$,
the {\bf base $b$ fixed-width $w$ expansion of $n$}  is
\[
(a_{w-1} \cdots a_1 a_0)_{b,w}
\]
where  $a_0, a_1, \ldots, a_{w-1}$ are nonnegative integers less than $b$ and
\[
n =  a_{w-1} b^{w-1} + \cdots + a_1b + a_0
\]
For $b$ an integer greater than $1$, $w$ a positive integer, $w'$ a positive  integer, and $x$ a real number
the {\bf base $b$ fixed-width expansion of $x$ with integer part width $w$  and fractional part width $w'$} is
\[
(a_{w-1} \cdots a_1 a_0 .  c_{1} \cdots c_{w'})_{b,w,w'}
\]
where  $a_0, a_1, \ldots, a_{w-1}, c_1, \ldots, c_{w'}$ are nonnegative integers less than $b$ and
$$x \geq a_{w-1} b^{w-1} +  \cdots + a_1 b + a_0 +  c_{1} b^{-1} + \cdots +  c_{w'} b^{-w'}$$
and
$$x < a_{w-1} b^{w-1} +  \cdots + a_1 b + a_0 +  c_{1} b^{-1} + \cdots +  (c_{w'} +1) b^{-w'}$$


For each set below, determine if it is empty, nonempty and finite, countably infinite, or uncountable.

{\it Challenge - not to hand in}: how would you prove this?

\begin{enumerate}
\item $B$
\item $S$
\item $\{ x \in \mathbb{N} ~\mid~ x = (4102)_3 \}$
\item $\{ x \in \mathbb{R} ~\mid~ \text{$x$ has a binary fixed-width $5$ expansion} \}$
\item $\{ x \in \mathbb{R} ~\mid~ x = (0.10)_{(2,1,2)} \}$
\end{enumerate} \end{enumerate}
\end{document}